\problem{6.8}{Hermiticity and spectrum of parity}
\textbf{(a)} In one dimension
$(\PP g)(x)=g(-x)$, and
\[
\braket f{\PP g}=\int_{-\infty}^\infty
f^*(x)g(-x)\dd x.
\]
Set $u=-x$, reverse integration limits:
\[
\braket f{\PP g}=
\int_{-\infty}^\infty f^*(-u)g(u)\dd u
=\braket{\PP f}g.
\]
Therefore $\boxed{\PP^\dagger=\PP}$.
The same coordinate substitution works in
three dimensions because the Jacobian's
absolute value equals one.
\textbf{(b)} Inverting coordinates twice gives
$\PP^2=I$. If $\PP\psi=\lambda\psi$,
then $\PP^2\psi=\lambda^2\psi=\psi$.
For a nonzero eigenstate
\begin{equation}
\boxed{\lambda^2=1,\qquad\lambda=\pm1.}
\end{equation}
Moreover $\PP^\dagger\PP=\PP^2=I$,
so parity is both Hermitian and unitary.

\problem{6.9}{Commutator and anticommutator parity classification}
For an operator $A$, let
$A'=\PP^\dagger A\PP=\PP A\PP$.
If $A$ is \emph{even}, $A'=A$.
Multiply on the left by $\PP$
and use $\PP^2=I$: $A\PP=\PP A$.
Hence
\[
\boxed{[\PP,A]=0\qquad\text{(even under parity)}.}
\]
If $A$ is \emph{odd}, $A'=-A$,
and the same manipulation gives
$A\PP=-\PP A$. Thus
\[
\boxed{\{\PP,A\}\equiv\PP A+A\PP=0
\qquad\text{(odd under parity)}.}
\]
Consequently a true scalar and
a pseudovector commute with parity;
a pseudoscalar and a true vector
anticommute with parity.
\phys{The scalar/vector distinction
concerns rotations; the true/pseudo
distinction concerns inversion.}

\problem{6.10}{Position and momentum are parity-odd}
Apply the triple product $\PP \x\PP$
to an arbitrary $f$:
\begin{align}
(\PP\x\PP f)(x)
&=\PP[x f(-x)]\\
&=-x f(x)=-(\x f)(x).
\end{align}
Therefore $\boxed{\PP\x\PP=-\x}$.
For momentum, using $\p=-\ii\hb\partial_x$,
\begin{align}
(\PP\p\PP f)(x)
&=\PP\left[-\ii\hb\frac{\dd}{\dd x}
f(-x)\right]\\
&=\PP[\ii\hb f'(-x)]
=\ii\hb f'(x)=-(\p f)(x).
\end{align}
Thus $\boxed{\PP\p\PP=-\p}$.
In three dimensions, apply the argument
to each coordinate:
\begin{equation}
\boxed{\PP^\dagger\mathbf r\PP=-\mathbf r,\qquad
\PP^\dagger\mathbf p\PP=-\mathbf p.}
\end{equation}
Their cross product
$\mathbf L=\mathbf r\times\mathbf p$
is consequently parity \emph{even}.

\problem{6.11}{Selection rule for a pseudovector or true scalar}
Let the two states have definite parities
$\PP\ket a=\eta_a\ket a$ and
$\PP\ket b=\eta_b\ket b$, where
$\eta_a,\eta_b=\pm1$.
For angular momentum, $\PP\mathbf L\PP=\mathbf L$.
Insert the identity as parity operators
on both sides:
\[
\bra b\mathbf L\ket a=
\bra b\PP^\dagger\mathbf L\PP\ket a
=\eta_b\eta_a\bra b\mathbf L\ket a.
\]
If $\eta_b\eta_a=-1$, the only possible
matrix element is zero:
\begin{equation}
\boxed{\bra{n'\ell'm'}\mathbf L\ket{n\ell m}=0
\quad\text{if }\ell+\ell'\text{ is odd}.}
\end{equation}
The same \emph{parity} rule holds for
any pseudovector or true scalar.
For $\mathbf L$ an even stronger
\emph{rotational} rule applies:
$\ell'=\ell$, because
$[\mathbf L,L^2]=0$ and $\mathbf L$
acts inside a fixed $\ell$ irreducible
multiplet.
\caveat{Parity gives a necessary condition,
not a guarantee of a nonzero matrix element.}

\problem{6.12}{Parity acts trivially on a pure spin-half spinor}
Let $\PP_s$ denote the $2\times2$ matrix
that acts only on the spin factor.
Spin is a pseudovector, so
$[\PP_s,S_i]=0$ for $i=x,y,z$.
Write
$\PP_s=\begin{pmatrix}a&b\\c&d\end{pmatrix}$.
From $[\PP_s,S_z]=0$ we obtain
$b=c=0$, so $\PP_s=\mathrm{diag}(a,d)$.
Then $[\PP_s,S_x]=0$ forces $a=d$:
$\PP_s=aI_2$.
The parity identity $\PP_s^2=I_2$
and Hermiticity give $a=\pm1$.
Choose the conventional intrinsic
spin parity $+1$:
\begin{equation}
\boxed{\PP_s=I_2.}
\end{equation}
The full electron state can still have
odd spatial parity $(-1)^\ell$.
This convention does not claim that an
electron and its antiparticle have the
same intrinsic parity.

\problem{6.13}{Electric dipole of the $1s$ state and an $n=2$ superposition}
\textbf{(a)} For an electron
$\mathbf p_e=-e\mathbf r$,
a parity-odd vector.
The hydrogen $1s$ state has even parity,
so $\av{\mathbf p_e}=-
\av{\PP^\dagger\mathbf p_e\PP}
=-\av{\mathbf p_e}$.
Therefore
$\boxed{\av{\mathbf p_e}_{1s}=0}$,
with no integration.
\textbf{(b)} The $n=2$ shell is
degenerate between $2s$ and $2p$.
Take the \emph{normalized} state
\begin{equation}
\boxed{\ket\psi=
(\ket{200}+\ket{210})/\sqrt2.}
\end{equation}
Its $x$ and $y$ dipoles vanish by
azimuthal integration. For $z=r\cos\theta$,
the two diagonal terms vanish by parity,
and only the cross term survives:
\[
\av z=\frac12
(\bra{200}z\ket{210}
+\bra{210}z\ket{200}).
\]
Use the real hydrogen functions
\[
R_{20}=\frac{(2-r/a_0)\ee^{-r/(2a_0)}}{
2\sqrt2\,a_0^{3/2}},\quad
R_{21}=\frac{(r/a_0)\ee^{-r/(2a_0)}}{
2\sqrt6\,a_0^{3/2}},
\]
and $Y_0^0=1/\sqrt{4\pi}$,
$Y_1^0=\sqrt{3/(4\pi)}\cos\theta$.
The angular integral is
$\int (Y_0^0)^*\cos\theta Y_1^0\dd\Omega
=1/\sqrt3$. The radial one is
\begin{align}
\int_0^\infty r^3R_{20}R_{21}\dd r
&=\frac{a_0}{8\sqrt3}
\int_0^\infty t^4(2-t)\ee^{-t}\dd t\\
&=\frac{a_0}{8\sqrt3}(2\cdot4!-5!)
=-3\sqrt3\,a_0.
\end{align}
Hence $\bra{200}z\ket{210}=-3a_0$,
and
\begin{equation}
\boxed{\av{\mathbf r}=-3a_0\hat{\mathbf z},
\qquad\av{\mathbf p_e}=+3ea_0\hat{\mathbf z}.}
\end{equation}
\phys{A degenerate superposition of
opposite-parity orbitals can carry a
static dipole moment even though each
pure orbital has zero dipole.}

\problem{6.14}{Finite rotations from infinitesimal planar matrices}
Put $\delta=\varphi/N$ and
\[
M_N=\begin{pmatrix}1&-\delta\\\delta&1\end{pmatrix}
=I+\delta J,\qquad
J=\begin{pmatrix}0&-1\\1&0\end{pmatrix}.
\]
\textbf{(a)} Since $J$ has eigenvectors
$v_-=(1,\ii)^T/\sqrt2$ and
$v_+=(1,-\ii)^T/\sqrt2$
with eigenvalues $-\ii$ and $+\ii$,
the unitary $S=(v_-\ v_+)$ gives
\[
\boxed{S^\dagger M_NS=
\operatorname{diag}(1-\ii\delta,1+\ii\delta).}
\]
\textbf{(b)} For each eigenvalue,
$\lim_{N\to\infty}
(1\mp\ii\varphi/N)^N=\ee^{\mp\ii\varphi}$.
Thus
\[
\lim_{N\to\infty}M_N^N=
S\,\operatorname{diag}(\ee^{-\ii\varphi},
\ee^{+\ii\varphi})S^\dagger.
\]
\textbf{(c)} Evaluate the product:
\begin{equation}
\boxed{\lim_{N\to\infty}M_N^N=
\ee^{\varphi J}=
\begin{pmatrix}
\cos\varphi&-\sin\varphi\\
\sin\varphi&\cos\varphi
\end{pmatrix}.}
\end{equation}
This exactly connects the
infinitesimal transformation
of Equation 6.30 with its finite
rotation in Equation 6.31.

\problem{6.15}{Rotational invariance of a scalar operator}
A scalar operator satisfies
$[L_i,f]=0$ for each Cartesian $i$.
Hence $[\mathbf n\cdot\mathbf L,f]=0$
for every unit $\mathbf n$.
Because
$\RR_{\mathbf n}(\varphi)
=\exp(-\ii\varphi\,\mathbf n\cdot\mathbf L/\hb)$,
every power in its exponential commutes
with $f$ and so does the complete rotation:
$\RR f=f\RR$.
Consequently
\begin{equation}
\boxed{\RR^\dagger f\RR
=\RR^\dagger\RR f=f.}
\end{equation}
The result holds for both true scalars
and pseudoscalars under \emph{proper}
rotations. Parity is a separate question.

\problem{6.16}{Infinitesimal rotation about the $y$ axis}
For a vector operator $\mathbf V$,
$[L_i,V_j]=\ii\hb\epsilon_{ijk}V_k$.
Under a small rotation of angle
$\delta$ about $y$,
\begin{align}
V_j'&=\RR_y^\dagger(\delta)V_j\RR_y(\delta)\\
&=V_j+\frac{\ii\delta}{\hb}[L_y,V_j]
+O(\delta^2)\\
&=V_j-\delta\epsilon_{yjk}V_k+O(\delta^2).
\end{align}
Using $\epsilon_{yxz}=-1$,
$\epsilon_{yzx}=+1$:
\[
V_x'=V_x+\delta V_z,\quad
V_y'=V_y,\quad
V_z'=V_z-\delta V_x.
\]
Thus
\begin{equation}
\boxed{\begin{pmatrix}V_x'\\V_y'\\V_z'\end{pmatrix}
=\begin{pmatrix}
1&0&\delta\\
0&1&0\\
-\delta&0&1
\end{pmatrix}
\begin{pmatrix}V_x\\V_y\\V_z\end{pmatrix}
+O(\delta^2).}
\end{equation}
This is the linear approximation to
the conventional $y$-axis rotation
matrix.

\problem{6.17}{Infinitesimal rotation of $\ket{n\ell m}$}
Expand the rotation operator to first order:
\[
\RR_{\mathbf n}(\delta)=
I-\frac{\ii\delta}{\hb}
(n_xL_x+n_yL_y+n_zL_z)+O(\delta^2).
\]
Use $L_x=(L_++L_-)/2$,
$L_y=(L_+-L_-)/(2\ii)$, giving
\[
\mathbf n\cdot\mathbf L=
n_zL_z+\frac{n_x-\ii n_y}{2}L_+
+\frac{n_x+\ii n_y}{2}L_-.
\]
On the angular-momentum basis,
\[
L_\pm\ket{n\ell m}=
\hb\sqrt{(\ell\mp m)(\ell\pm m+1)}
\ket{n\ell,m\pm1},
\quad L_z\ket{n\ell m}=m\hb\ket{n\ell m}.
\]
Therefore
\begin{align}
\RR_{\mathbf n}(\delta)\ket{n\ell m}
&=\sum_{m'=-\ell}^{\ell}
D_{m'm}^{(\ell)}\ket{n\ell m'}+O(\delta^2),\\
\boxed{D_{m'm}^{(\ell)}}&=
\delta_{m'm}-\ii\delta\bigg[
n_zm\,\delta_{m'm}\\
&\quad+\frac{n_x-\ii n_y}{2}
\sqrt{(\ell-m)(\ell+m+1)}\,
\delta_{m',m+1}\\
&\quad+\frac{n_x+\ii n_y}{2}
\sqrt{(\ell+m)(\ell-m+1)}\,
\delta_{m',m-1}
\bigg].
\end{align}
The principal $n$ and total orbital
$\ell$ labels remain fixed, while
the projection $m$ mixes with its
neighbors. As $D=I-\ii\delta A$
with $A=A^\dagger$, $D^\dagger D
=I+O(\delta^2)$, as required
by unitarity.
